\section{Guidance and Open Problems}
\label{sec:discussion}

Previous research has found that decision making in AI-assisted contexts improves when humans and AI complement each other, taking their respective weaknesses into account \cite{bansal2019beyond}. Without an understanding of AI weaknesses, humans may over- or under-rely on AI for assistance \cite{schoeffer2024ai}. On the human side, three well-established weaknesses or challenges for which AI can \textit{potentially} compensate are: 1. The challenge of keeping up with security guidance and tools, 2. User mental models of security that may not be aligned with good security practices \cite{wash2011influencing} and 3. Barriers to user acceptance of security advice \cite{redmiles2016learned}.  We consider the findings of~\Cref{findings}  in the context of these challenges to distill guidance for end-users and model developers toward improving human-LLM interaction in the user security assistance context. We also identify places where additional research is needed to extend this guidance.

\subsection{User Guidance}

We discuss what our evaluation suggests regarding security topics and question types best suited to LLMs, how user questions are expressed (i.e., prompt engineering), and how users can assess LLM responses. Note that while we report performance numbers for GPT only, all the errors discussed in this section were also evident with Llama and Gemini, except where otherwise indicated.\footnote{Our analysis of Llama and Gemini responses focused on questions that GPT responded to incorrectly, so while we know similar errors occur with Llama and Gemini, we do not know how they performed overall.}

\paragraph{When to use LLMs.} Based on our evaluation, we discourage users from relying on LLMs for procedural questions (e.g., ``How do I…?’’), regardless of security topic. Our evaluation found a high rate of errors (31\% for GPT) and omissions of important information (43\% for GPT) in response to procedural questions.

Compatible with research in other domains, LLMs struggled with questions requiring reasoning or comparison, such as security trends. For example, when asked ``What age groups most report fraud or scams?'', GPT calls out 4 age groups, including teens, adults and seniors, when in fact seniors are less likely to report fraud and scams than younger people \cite{ftc_scams_all_ages_2022}.

Another challenging area is questions that are sensitive to recent events, such as queries regarding trust and reputation. For example, when asked ``Can I get a free password manager?'', GPT suggests Lastpass (among others), a password manager that is no longer recommended due to recent breaches.

Finally, we encourage care around questions in the areas of authentication and antivirus, particularly when asking product, platform, or company-related questions, since significant weakness in response quality was found in both. GPT errors exceeded 30\% in both categories, and the same weakness was present in both Llama and Gemini responses. 

Areas of strength for LLMs are definitions of security concepts (e.g., ``What is malware?''), basic risk assessment questions (e.g., ``Will 2 factor authentication protect me?'') and core functionality of established products (e.g., ``How does Keeper password manager work?''). This is compatible with the fact that factual and conceptual questions are the most likely to be answered correctly (\Cref{tab:gpt-evaluation-dist-by-topics}). 

\paragraph{How to communicate with LLMs.} We used questions from published FAQs and question templates that our review found to be written in ``plain’’ English (i.e., no technical, rare or otherwise lesser-known words and phrases \cite{williams2004legal}). We conducted a small experiment to test the sensitivity of all three LLMs to rephrased questions---which are organic user queries seen on Google search---from this corpus (Section~\ref{sec:reproducibility}) and did not find sensitivity to phrasing.


In addition, several other open prompt engineering questions would allow for important extensions of user guidance in this area. For example, we found a substantial number of indirect responses (17\% with GPT) that delay information delivery and could cause user attention to shift elsewhere. We term a portion of these responses ``LLM-splaining’’ because they provide information that is likely already known by the user. The unnecessary repetition of information may diminish user trust \cite{zeffane2011communication}. For example, to the question, ``How do I set up Okta Verify?’’, GPT responds with 2 sentences explaining Okta Verify before beginning to answer the user question: ``Okta Verify is an authentication app that works in conjunction with Okta, a secure identity management service. It provides a secure way…’’. Prompts that encourage LLMs to consistently address user questions directly are an open problem.

Finally, we note that while our questions did not contain any sensitive information, security questions can involve sensitive data so user prompts may be sensitive, particularly in light of patterns of over-sharing with LLMs that have been found \cite{zhang2024s}. The relationship between LLM response quality and the inclusion of sensitive information is largely unexplored and would be a useful addition to user guidance.


\paragraph{User assessment of LLM responses.} While this may change with LLM-enabled search, none of the LLMs we experimented with routinely cite sources when answering end-user security questions. Indeed, URLs are present in less than 5\% of the GPT answers in our evaluation. We recommend users explicitly ask for authoritative sources supporting LLM responses, and consider requesting that LLMs find supporting data from government or nonprofit sources of the user’s choosing (e.g., NIST, FTC).

As discussed in \cite{barman2024beyond}, ``simple heuristics to detect potential
errors or misinformation from LLMs’’ should be part of a user toolkit for safe interaction with LLMs. Towards such a toolkit we see evidence of two heuristics in our data given that errors in some of the evaluation criteria are predictive of other errors. In particular, indirect GPT responses ($n=153$) have errors in accuracy, thoroughness and relevance with probability greater than $.58$. Similarly, irrelevant GPT answers ($n=15$) have errors in accuracy, thoroughness or directness with probably greater than $.73$. The support for each of these heuristics is relatively low, so it is an open problem to confirm that these heuristics hold. That said, they suggest a promising direction for helping users safely interact with LLMs since indirect responses and irrelevant content may be more easily recognized than the other criteria and serve as an early user warning that a response may be low quality. While there has been recent progress in computational approaches to predicting hallucinations \cite{farquhar2024detecting}, \textit{user-friendly} heuristics are an important open problem.

\add{We note that our study %
informs user guidance for interacting with LLMs but does not explore user perception of LLM responses. Understanding the applicability and understandability of LLM responses to user security questions is an open problem.}

\subsection{Developer Guidance}
In this section, we discuss the primary areas of performance problems in our evaluation, organized according to the user challenges described at the beginning of Section~\ref{sec:discussion}. While we suggest directions for improvements, we are not prescriptive, as model developers are better positioned to identify root causes and craft solutions.

\paragraph{Generating correct security content.} Our evaluation suggests augmenting LLM training data with more security research and reducing the support of security product marketing materials will improve LLM performance. For example, all LLMs in our evaluation routinely provide stale password guidance despite ample research attesting to the usability shortcomings of such guidance. In addition, the privacy and security risks of default app permissions, a topic well-covered by research \cite{felt-permissions-2011, chia-permission-risks-2012, calciati-auto-permissions-2020}, are often overlooked in LLM responses. Finally, the incorrectly inflated capabilities of certain security technologies (e.g., HTTPS) suggest marketing materials may have high support in LLM training data. For example, “Authy takes security very seriously", and “Dashlane is a well-established and reputable password manager that has earned the trust of millions of users worldwide." Overly favorable product content may be perceived as marketing information and can be devalued by users \cite{redmiles2016learned}. 



Questions involving specific companies, products, or platforms are often challenging for LLMs. For example, when no company, product, or platform is mentioned, more than 60\% of GPT responses are perfectly accurate, thorough, and relevant. In contrast, this number drops to less than 44\% when a specific company, product, or platform is part of the prompt. We suggest that the presence of a company name, product, or platform may be useful for developers in determining when to prioritize authoritative sources.


Finally, while few in number, our evaluation did surface 10 examples in which the human-LLM interaction completely failed because the LLM was not responsive to the question. In 7 of these (6 for Gemini and 1 for Llama), this appeared to be due to mistriggered guardrails; the 3 others were GPT misinterpretations (e.g., to the question, ``Can you write your login credentials down, really?’’ GPT responds, ``I'm sorry for any confusion, but as an AI, I don't possess personal login credentials…”). We suggest that developers consider how to better confirm the user intent before delivering a non-answer.

In short, while those with a complete view of the model development process are best positioned to determine how to address these errors as well as the hallucinations and omitted risks discussed in Section~\ref{findings}, we emphasize that answers to questions in these areas are well documented in authoritative security sources (e.g., peer-reviewed research papers, government reports, product and platform documentation, etc.) and so it appears models either do not have sufficient access to this data or aren’t prioritizing data sources appropriately.

\paragraph{Aligning mental models with sound security practices.} Our question corpus was gathered to broadly represent user concerns in security, rather than focus on areas known to have potentially problematic user models. That said, we do see evidence of known user mental models in LLM answers. For example, \cite{wash2010folk} found that the term ``hacker'' is often used to broadly describe bad actors in a security context regardless of the actor's attack strategy or relationship to the user. We see a similar broad use of the term in LLM responses, GPT characterizes actors launching technical attacks as hackers (``The longer the password, the more combinations a hacker has to try to crack it.") as well as scammers (``Hackers can use sophisticated phishing schemes to trick users"). In addition, the inflation of the protective power of HTTPS noted in Section~\ref{findings} was also found in \cite{herbert2023world} among participants in India.

We emphasize that LLM perpetuation of user mental models is only problematic if those models are associated with risky security practices. For example, ``hacker'' terminology may suppress the fact that security risks also exist with people a user knows (e.g., intimate partner or elder abuse).

Finally, we note while some LLM error patterns may not specifically perpetuate problematic mental models, they may still contribute to areas of confusion. While we are unaware of previous research around the mental models of linking VPNs to phishing protection, the protective power of VPNs is a well-established area of user confusion (e.g., \cite{dutkowska2022and, binkhorst2022security}).

We encourage model developers to use security user mental models research to explore potential areas of bias and consider relying more heavily on authoritative sources in those areas.
 
\paragraph{Encouraging adoption of security information.} We recommend developers tailor the balance between content and sources based on the nature of the user question. In particular, consider deferring to sources in the cases of procedural and UI-dependent questions,  or when contact information or URLs are part of the response. For instance, GPT provided email addresses in 6 responses, none of which were valid at the time of our evaluation. Also, in 3 responses, GPT provided 5 (out of 123 URLs across 47 responses) inaccessible URLs. Our evaluation indicates making strategic choices between the balance of content and sources in LLM answers will improve response quality and may increase user acceptance; there is evidence that acceptance can be sensitive to the authoritativeness of the source, for example in areas that are less familiar to the user \cite{redmiles2016learned}. 

Given that user security questions may be time-sensitive, we recommend LLMs use a direct communication style in response to such questions. In addition, developers should consider signaling to the user when they are not able to deliver a high-confidence response, similar to the warnings many LLMs already provide about potential limitations due to the date range of training data. Finally, we raise the question of whether explanations can also help the end user better detect errors in the end user security context \cite{bansal2021does}.






